\unitopen{03}
\begin{candos}
\item liczyć po niemiecku do stu i zrozumieć liczbę usłyszaną w rozmowie;
\item podać swój wiek, numer telefonu i adres e-mail;
\item przeliterować nazwisko i adres e-mail przez telefon;
\item wypełnić prosty formularz (imię, nazwisko, adres, kod pocztowy).
\end{candos}

\section{Słownictwo}
\vocabtable{03}

\section{Liczby 0--100}
Liczby 13--19 powstają z dodania \de{-zehn} do jedności: \de{dreizehn}, \de{vierzehn}, \de{fünfzehn},
\de{sechzehn}, \de{siebzehn}, \de{achtzehn}, \de{neunzehn}. Uwaga na \de{sechzehn} i \de{siebzehn}: gubią końcówkę
(\de{sechs} $\rightarrow$ \de{sech-}, \de{sieben} $\rightarrow$ \de{sieb-}). Dziesiątki kończą się na \de{-zig}
(\de{dreißig}, nie \de{dreizig}).

\begin{zasada}[Jedności przed dziesiątkami]
Od 21 do 99 mówisz najpierw jedności, potem dziesiątki, połączone słowem \de{und}, pisane razem:
\de{21 = einundzwanzig}, \de{34 = vierunddreißig}, \de{67 = siebenundsechzig}, \de{99 = neunundneunzig}.
W polskim jest odwrotnie. Przez telefon zapisuj liczbę, którą słyszysz, \emph{od końca}: najpierw ostatnia cyfra.
\end{zasada}

\begin{center}\small
\begin{tabular}{@{}rlrlrl@{}}
13 & \de{dreizehn} & 17 & \de{siebzehn} & 40 & \de{vierzig}\\
14 & \de{vierzehn} & 18 & \de{achtzehn} & 60 & \de{sechzig}\\
15 & \de{fünfzehn} & 19 & \de{neunzehn} & 70 & \de{siebzig}\\
16 & \de{sechzehn} & 20 & \de{zwanzig} & 100 & \de{(ein)hundert}\\
\end{tabular}
\end{center}

\section{Wiek, telefon, e-mail}
Wiek: \de{Ich bin 30 Jahre alt.} Pytanie: \de{Wie alt bist du?} / \de{Wie alt sind Sie?} Numer telefonu czyta
się parami lub pojedynczo: \de{+41 44 123 45 67} = \de{plus vier-eins, vier-vier, eins-zwo-drei, vier-fünf,
sechs-sieben}. E-mail: \de{max.muster@beispiel.ch} = \de{Max Punkt Muster ät Beispiel Punkt Ce-Ha}.
Znaki: \de{Punkt} (.), \de{Bindestrich} (-), \de{Unterstrich} (\_), \de{ät} (@).

\begin{falle}
Po polsku mówisz „mam 30 lat”, po niemiecku \de{Ich \emph{bin} 30 Jahre alt} (jestem, nie mam). Nie używaj
\de{haben} do wieku (\de{Ich habe 30 Jahre} jest błędem). I drugie: w niemieckim dokumencie
liczbę z tysiącami zapisujesz z kropką (\de{1.000}) albo spacją, ułamek z przecinkiem (\de{3,5}), odwrotnie niż
w angielskim.
\end{falle}

\begin{szwajcaria}
W Szwajcarii zapiszesz numer jako \de{044 123 45 67} (lokalnie) albo \de{+41 44 123 45 67}.
Separator tysięcy często jest apostrofem: \de{1'000'000}. Kwoty: \de{CHF 1'250.50} (z kropką dziesiętną).
Piszesz \de{dreissig}, nie \de{dreißig}: szwajcarski standard nie używa \de{ß}.
\end{szwajcaria}

\begin{praca}
Rekruter przez telefon: \de{Können Sie bitte Ihren Namen buchstabieren?} Odpowiadasz literowaniem
(\de{Ka-O-En-Er-A-De}) i podajesz e-mail wolno, znak po znaku. Do literowania ćwicz zwłaszcza
\de{E}, \de{I}, \de{G}, \de{J}, \de{V}, \de{W}.
\end{praca}

\section{Dialogi}
\textbf{Dialog 1 --- telefon od rekrutera}
\begin{dialog}
\turn{A}{Guten Tag, hier ist Anna Keller von der Firma Alpha. Spreche ich mit Herrn Nowak?}{Dzień dobry, mówi Anna Keller z firmy Alpha. Czy rozmawiam z panem Nowakiem?}
\turn{B}{Ja, am Apparat. Guten Tag, Frau Keller.}{Tak, przy telefonie. Dzień dobry, pani Keller.}
\turn{A}{Wie ist Ihre E-Mail-Adresse?}{Jaki jest Pana adres e-mail?}
\turn{B}{piotr Punkt nowak ät beispiel Punkt com.}{piotr kropka nowak małpa beispiel kropka com.}
\turn{A}{Und Ihre Telefonnummer?}{A numer telefonu?}
\turn{B}{Null-sieben-neun, vier-fünf-zwo, eins-drei, sechs-sieben.}{Zero-siedem-dziewięć, cztery-pięć-dwa, jeden-trzy, sześć-siedem.}
\end{dialog}

\section{Ćwiczenia}
\begin{cwiczenia}
\exitem{Napisz słownie: 13.}{\de{dreizehn}}
\exitem{Napisz słownie: 21.}{\de{einundzwanzig}}
\exitem{Napisz słownie: 34.}{\de{vierunddreißig}}
\exitem{Napisz słownie: 67.}{\de{siebenundsechzig}}
\exitem{Napisz słownie: 99.}{\de{neunundneunzig}}
\exitem{Policz: \de{Was ist dreißig plus zwölf?}}{\de{zweiundvierzig}}
\exitem{Policz: \de{Was ist hundert minus fünfundzwanzig?}}{\de{fünfundsiebzig}}
\exitem{Powiedz, ile masz lat (np.\ 28).}{\de{Ich bin achtundzwanzig Jahre alt.}}
\exitem{Zadaj pytanie o wiek koledze (nieformalnie).}{\de{Wie alt bist du?}}
\exitem{Przeczytaj adres \de{max.muster@beispiel.ch}.}{\de{Max Punkt Muster ät Beispiel Punkt Ce-Ha}}
\exitem{Co jest błędne w: \de{Ich habe dreißig Jahre.}?}{Wiek podaje się z \de{sein}: \de{Ich bin dreißig Jahre alt.}}
\exitem{Przetłumacz: „Jaki jest Pana numer telefonu?”}{\de{Wie ist Ihre Telefonnummer?}}
\end{cwiczenia}

\weekend{Przygotuj \emph{wizytówkę}: imię, nazwisko, adres, kod pocztowy, telefon, e-mail --- po niemiecku
i na głos (z literowaniem). Poproś asystenta AI, żeby zadał Ci pytania z formularza (\de{Wie ist Ihr Nachname?}
\de{Wie ist Ihre Postleitzahl?}) i poprawił odpowiedzi.}
